\section{Introduction}\label{sec:intro}

Small multirotor drones are becoming part of the infrastructure of smart cities. Municipal operators fly them to
inspect bridges and power lines, to search for missing persons, to survey traffic and crowds after an incident and to
deliver medical supplies, and they increasingly fly them as fleets dispatched from shared depots rather than as single
aircraft under a pilot's hand~\cite{mohamed2020uav,gharibi2016internet,prevot2016utm}. A fleet operator takes the same
few decisions for every sortie: whether a drone may launch for a task, at which altitude and along which path it
returns when its battery or the operator calls it home, and how high and how densely a scan must be flown so that the
camera can still recognise what it is looking for.

These decisions are \emph{coupled}. The right return altitude depends on the buildings along the return line; when the
return must start depends on the energy needed to fly that path against the wind; the right scan height depends on the
visibility, which in turn sets the number of strips, the flight time per sortie and thus the number of drones. A
decision that sees only one of these factors is not merely less accurate; it is wrong in a systematic way that depends
on the city and on the weather. Yet the tools of today keep the factors apart. Flight stacks return at a fixed
altitude above the take-off point~\cite{meier2015px4}; drone simulators model a vehicle in a uniform
wind~\cite{shah2017airsim,koenig2004gazebo}; urban digital twins display real geometry but do not
simulate~\cite{batty2018digitaltwins,ketzler2020digitaltwins}; and weather enters operations, if at all, as a go/no-go
rule of thumb~\cite{gao2021weather}. Each module therefore decides with its own partial model, and two modules that
need the same quantity may compute it differently.

We measure what this costs in six real cities and one synthetic city (Sect.~\ref{sec:bg:study}). A return at the default altitude of
PX4~\cite{meier2015px4}, a widely used open-source flight stack, enters the building clearance envelope on 45\% of return lines, and no fixed altitude is both safe and
cheap in any of the cities. A return trigger with its own headwind model aborts up to a quarter of legs that the
planner's energy model considers feasible. A scan plan made for clear air captures 5\% of the targets in fog while
predicting full coverage. And the obvious fix, querying geometry, wind and energy for every drone at every tick, costs
more than thirteen CPU cores for a fleet of 1,000 drones. Coupled decisions therefore need two things at once: one
consistent model of the city, the weather and the vehicle that every module queries, and a runtime that makes those
queries cheap enough for a whole fleet.

This paper presents \emph{Drone4D}, a space-time (4D) world runtime for urban drone fleets that provides both, and a
decision layer built on it. The runtime serves one model of city geometry, of a weather field that varies in space and
time, and of vehicle energy and perception, to the planner, to the simulated fleet and to the operator's browser
(Sect.~\ref{sec:design}). Its design follows one rule: \emph{a quantity that two modules need is computed by one
function}. A versioned World Package turns city point clouds into a geometry service with exact and bounded corridor
queries; an analytic environment field gives wind, visibility and air density at any point and time and is evaluated
identically on the server and in the browser; the fleet runtime advances up to 1,000 drones on one simulation clock with
staged, batched queries and refreshes return plans on a rotating schedule; and a budgeted point-cloud client streams
the city to ordinary browsers without a server GPU. On top of the runtime, the decision layer couples the launch
precheck, return planning, a perception-driven scan planner and capability-based task delegation (contract net) to the
shared models. Switches that hide individual inputs from a decision, while the simulated ground truth stays fully
coupled, turn every baseline and ablation into a configuration of the same system.

Our evaluation in closed-loop campaigns over seven cities and twelve weather conditions (Sect.~\ref{sec:eval}) shows
what the coupling enables. Compared with a PX4-style return, the coupled policy reduces unsafe sorties from
\PXUNSAFE\% to \CPUNSAFE\% and completes \CPDONE\% of sorties instead of \PXDONE\%. When a weather front is forecast, a
precheck that queries the field at the times the legs will be flown cuts unsafe sorties from 40--89\% to 4--11\%. The
scan planner raises recognition capture in fog from 5\% to 89\% and states the price in drones before launch, and
delegation on the shared models raises request success from 42\% to 59\% without stranding a drone. The coupling
costs little: on one CPU core the runtime advances 1,000 drones at \STAGERTF{} times real time, and keeping every return
plan current costs 1.25 core-milliseconds per simulated second.

In summary, this paper makes the following contributions:
\begin{itemize}
\item A data-driven study in six real cities and one synthetic city that quantifies how decisions based on partial
models fail: returns that fly into buildings, modules that contradict each other in wind, scans that silently miss their
targets, and the cost of naive coupling (Sect.~\ref{sec:background}).
\item The design and implementation of Drone4D, a space-time world runtime that serves one consistent model of geometry, weather and
energy to the planner, the fleet runtime and the browser, with batched queries and a rotating refresh that keep coupled
decisions within the real-time budget of a 1,000-drone fleet on a CPU-only server (Sect.~\ref{sec:design}).
\item A coupled decision layer, comprising a launch precheck with a space-time variant, a geometry- and energy-aware
return, a perception-driven scan planner and capability-based task delegation, together with decision-visibility
switches that make every baseline and ablation a configuration of one system (Sect.~\ref{sec:design:decisions}).
\item An evaluation over 60,480 closed-loop sorties, weather fronts, scan and delegation campaigns, micro-benchmarks,
cross-device streaming and consistency tests, which quantifies the gain of each coupling and its cost
(Sect.~\ref{sec:eval}).
\end{itemize}

The remainder of the paper is organised as follows. Section~\ref{sec:background} presents the motivational study, the
resulting challenges and related work. Section~\ref{sec:design} describes the runtime and the decision layer,
Sect.~\ref{sec:eval} evaluates them, and Sect.~\ref{sec:conclusion} concludes.
